\clearpage
\VerdierPhysicalPageAnchor{0155}%
\begingroup\small
\setlength{\parskip}{0pt}\setlength{\emergencystretch}{1.5em}

\VerdierSemanticAnchor{ch3-3-2-2-admissible}
$\mathcal D^\wedge$ 的对象若为两个可表示函子之间某个态射的像，就称为
\emph{可容许对象}。设
$\operatorname{Im}\colon\mathcal{F}\ell_1(\mathcal D)\to\mathcal D^\wedge$
为函子
\VerdierSemanticAnchor{ch3-3-2-2-1}
\[
\operatorname{Im}\bigl(X\xrightarrow fY\bigr)=h(f)\text{ 的像}.
\tag{3.2.2.1}
\]
$\operatorname{Im}$ 是加法函子，并把可忽略态射送到零态射。因此它唯一
分解为加法函子（仍滥用记号记为）
\VerdierSemanticAnchor{ch3-3-2-2-2}
\[
\operatorname{Im}\colon\mathrm A(\mathcal D)\longrightarrow\mathcal D^\wedge.
\tag{3.2.2.2}
\]
复合函子
$\operatorname{Im}\circ\mathcal H\colon
\mathcal D\longrightarrow\mathrm A(\mathcal D)
\longrightarrow\mathcal D^\wedge$
就是函子 $h$。$\mathcal D^\wedge$ 的对象同构于
$\mathrm A(\mathcal D)$ 中某个对象在 $\operatorname{Im}$ 下的像，当且仅当它可容许。

\VerdierSemanticAnchor{ch3-3-2-3}
\medskip\noindent\textbf{命题 3.2.3.}\quad
\textit{函子 $\operatorname{Im}\colon\mathrm A(\mathcal D)\to
\mathcal D^\wedge$ 充分忠实。}

\VerdierSemanticAnchor{ch3-3-2-3-morphism-square}
设
\[
\begin{tikzcd}[column sep=4em,row sep=2.4em]
X\arrow[r,"f"]\arrow[d,"g"']&Y\arrow[d,"g'"]\\
X'\arrow[r,"f'"']&Y'
\end{tikzcd}
\qquad=(g;g')
\]
为 $\mathcal{F}\ell_1(\mathcal D)$ 中的态射，并以
$\operatorname{Im}((g;g'))$ 表示其像。图式
\VerdierSemanticAnchor{ch3-3-2-3-canonical-diagram}
\[
\begin{tikzcd}[column sep=4em,row sep=2.2em]
X\arrow[r,"\operatorname{id}_X"]\arrow[d,"\operatorname{id}_X"']&X\arrow[d,"f"]\\
X\arrow[r,"f"]\arrow[d,"f"']&Y\arrow[d,"\operatorname{id}_Y"]\\
Y\arrow[r,"\operatorname{id}_Y"']&Y
\end{tikzcd}
\]
经 $\operatorname{Im}$ 给出
\[
h(X)\xrightarrow{\operatorname{Im}((\operatorname{id}_X;f))}
\operatorname{Im}(f)\xrightarrow{\operatorname{Im}((f;\operatorname{id}_Y))}h(Y).
\]
\VerdierSemanticAnchor{ch3-3-2-3-1}
并且
\[
\operatorname{Im}((f;\operatorname{id}_Y))\circ
\operatorname{Im}((\operatorname{id}_X;f))=h(f).
\tag{3.2.3.1}
\]
\endgroup
